
Large language models can explain code errors in detail, describing why an algorithm handles edge cases incorrectly or why a data structure choice leads to inefficiency. Yet these explanations are often incorrect. Analysis of Self-Refine by Madaan et al. (2023) indicates that 94\% of its failures trace to bad feedback: 33\% identify the wrong location and 61\% propose incorrect fixes. Meanwhile, execution-based feedback from unit tests provides ground-truth correctness signals but offers no insight into why code fails---only that it fails.

This work frames the tension as one between \textbf{signal fidelity} and \textbf{semantic richness}. Execution feedback achieves high fidelity (tests are deterministic) but low richness (no explanation). AI feedback achieves high richness (detailed explanations) but low fidelity (frequent errors). These dimensions are orthogonal: high fidelity does not preclude high richness.

The proposed mechanism, \textbf{Execution-Verified AI Feedback (EVAF)}, uses execution to verify AI feedback rather than replace it. Given code that fails tests, an AI model generates a critique and proposed fix. The fix is executed against unit tests. If all tests pass, the critique is accepted as verified feedback; otherwise, it is rejected.

The \textbf{accept rate}---the fraction of AI suggestions surviving verification---determines EVAF's viability. If accept rate is below 10\%, EVAF degenerates to pure execution feedback because AI is nearly always wrong. If accept rate exceeds 90\%, gating adds no value because AI is nearly always correct. The target range is 20-60\%, indicating selective filtering.

This work implements EVAF using CodeT5-770M and CodeLlama-7b-Instruct on HumanEval. The contribution is architectural: a working implementation of the EVAF pipeline. However, the experimental evaluation did not complete due to infrastructure failure during execution. No metrics were computed, and the hypothesis remains untested.
